% META: {"id": "1NSI-LANG-EVAL-B-corrige", "chapitre": "1NSI-LANGAGE", "type_objet": "corrige_evaluation", "evaluation_ref": "1NSI-LANG-EVAL-B", "capacites": ["P-LANG-01", "P-LANG-02", "P-LANG-03A", "P-LANG-03B", "P-LANG-03C", "P-LANG-04", "P-LANG-05"], "mode_creation": "ex_nihilo", "sources_inspiration": [], "status": "verified"}
% BEGIN-VERIFY
% def moyenne_ponderee(valeurs, poids):
%     assert len(valeurs) == len(poids)
%     assert sum(poids) > 0
%     total = 0
%     for i in range(len(valeurs)):
%         total = total + valeurs[i] * poids[i]
%     return total / sum(poids)
% assert moyenne_ponderee([8, 12, 16], [1, 1, 2]) == 13.0
% def produit_negatifs_bugue(liste):
%     prod = 0
%     for x in liste:
%         if x < 0:
%             prod *= x
%     return prod
% assert produit_negatifs_bugue([-2, -3, -4]) == 0
% import math
% assert math.gcd(60, 48) == 12
% def indice_max(nombres):
%     assert len(nombres) > 0, "la liste doit etre non vide"
%     meilleur = 0
%     for i in range(1, len(nombres)):
%         if nombres[i] > nombres[meilleur]:
%             meilleur = i
%     return meilleur
% for essai in ([3], [1, 5, 2], [4, 4, 1], [-3, -1, -7], [2, 9, 9, 1]):
%     k = indice_max(essai)
%     assert 0 <= k < len(essai)
%     assert all(valeur <= essai[k] for valeur in essai)
%     assert all(essai[j] < essai[k] for j in range(k))
% assert indice_max([4, 4, 1]) == 0
% assert indice_max([2, 9, 9, 1]) == 1
% try:
%     indice_max([])
% except AssertionError as e:
%     assert "non vide" in str(e)
% else:
%     raise AssertionError("la precondition n'a pas ete declenchee")
% END-VERIFY
\begin{corrige}{1NSI-LANG-EVAL-B-EX1}
\textbf{1.} Précondition : \lstinline{valeurs} et \lstinline{poids} ont la même longueur,
et la somme des poids est strictement positive.

\textbf{2.}
\begin{python}
def moyenne_ponderee(valeurs, poids):
    assert len(valeurs) == len(poids)
    assert sum(poids) > 0
    total = 0
    for i in range(len(valeurs)):
        total = total + valeurs[i] * poids[i]
    return total / sum(poids)
\end{python}

\textbf{3.} $\dfrac{8\times1+12\times1+16\times2}{1+1+2}=\dfrac{8+12+32}{4}=\dfrac{52}{4}=13$.
\end{corrige}

\begin{corrige}{1NSI-LANG-EVAL-B-EX2}
\textbf{1.} \lstinline{produit_negatifs_bugue([-2, -3, -4])} renvoie $0$. Le résultat
attendu est $-24$ ($(-2)\times(-3)\times(-4)$).

\textbf{2.} L'accumulateur \lstinline{prod} est initialisé à $0$ : dès la première
multiplication (\lstinline{prod *= x}), $0\times x=0$ quel que soit $x$, donc
\lstinline{prod} reste $0$ pour toujours. Il faut initialiser un accumulateur de
\textbf{produit} à $1$ (l'élément neutre de la multiplication), pas à $0$ :
\begin{python}
def produit_negatifs(liste):
    prod = 1
    for x in liste:
        if x < 0:
            prod *= x
    return prod
\end{python}

\textbf{3.} \lstinline{import math} puis \lstinline{math.gcd(60, 48)}, qui renvoie $12$.
\end{corrige}

\begin{corrige}{1NSI-LANG-EVAL-B-EX3}
\textbf{1.} Prototype : \lstinline{indice_max(nombres)} ; \lstinline{nombres} est une liste
de nombres ; la fonction renvoie un entier, l'indice du plus grand élément.
\textbf{Precondition :} la liste est non vide -- sans elle, il n'existe aucun indice a
renvoyer.

\textbf{2.} Deux postconditions verifiables :
\begin{itemize}
  \item l'indice renvoye est un indice valide : \lstinline{0 <= k < len(nombres)} ;
  \item aucune valeur ne depasse celle a cet indice, et aucun indice plus petit ne porte
  deja ce maximum -- ce qui exprime la règle « le plus petit indice en cas d'égalité ».
\end{itemize}

\begin{python}
def indice_max(nombres):
    assert len(nombres) > 0, "la liste doit etre non vide"
    meilleur = 0
    for i in range(1, len(nombres)):
        if nombres[i] > nombres[meilleur]:
            meilleur = i
    return meilleur


essai = [2, 9, 9, 1]
k = indice_max(essai)
assert 0 <= k < len(essai)
assert all(valeur <= essai[k] for valeur in essai)
assert all(essai[j] < essai[k] for j in range(k))
\end{python}

La comparaison stricte \lstinline{>} est ce qui garantit la seconde postcondition : avec
\lstinline{>=}, une égalité ferait avancer l'indice retenu.

\textbf{3.} \textbf{Trait commun :} les deux langages ont des fonctions nommees, des
variables, une boucle qui parcourt les indices, et la même logique de comparaison --
l'algorithme s'ecrit a l'identique. \textbf{Trait particulier a Python :} les blocs sont
delimites par l'\textbf{indentation}, et \lstinline{range} engendre les indices.
\textbf{Trait particulier a JavaScript :} les blocs sont delimites par des \textbf{accolades},
et la boucle \lstinline{for} declare explicitement son initialisation, sa condition d'arret
et son increment.
\end{corrige}
